% Fragmento para Overleaf: sustituye \chapter{Marco Teorico} ... hasta antes de \chapter{Cronograma}.
% Hacer copia de main.tex. Compilar pdfLaTeX + biber.
% Texto alineado con docs/academico/01_estado_del_arte.md y 02_metodologia.md.

\chapter{Estado del arte}

Este capítulo no es un manual de dinámica de gases. Sitúa el proyecto ---un \textit{framework} compacto para las ecuaciones de Euler por volúmenes finitos--- frente a los métodos que el artefacto debe comparar, al software científico existente y a las prácticas de ingeniería de software científico. No se declara una revisión sistemática protocolizada.

\section{Fundamentos numéricos pertinentes}

Las ecuaciones de Euler para un gas ideal no viscoso y sin conducción de calor son el prototipo de sistema hiperbólico conservativo sobre el que se diseñan esquemas de captura de choques \parencite{toro2009,leveque2002,hirsch2007}. El vector de estado $U=(\rho,\rho\mathbf{v},E)$ y el flujo $F(U)$ permiten conservar masa, momento y energía en la malla, condición necesaria para localizar discontinuidades \parencite{toro2009}. Las soluciones pueden incluir choques, rarefacciones y discontinuidades de contacto; el contacto (salto de densidad con presión y velocidad continuas) es el criterio que distingue la familia HLL de HLLC \parencite{toro2009}.

El método de volúmenes finitos aproxima el flujo en cada cara mediante un problema de Riemann local \parencite{leveque2002}. En una malla uniforme unidimensional,
\begin{equation}
\frac{\mathrm{d}U_i}{\mathrm{d}t}=-\frac{1}{\Delta x}\bigl(\hat{F}_{i+1/2}-\hat{F}_{i-1/2}\bigr).
\end{equation}
MUSCL eleva el orden en zonas suaves con una reconstrucción lineal limitada; el paso explícito queda acotado por CFL \parencite{toro2009,hirsch2007}. El solver de Riemann exacto sirve como \textbf{referencia de verificación}, no como dato de laboratorio \parencite{toro2009,oberkampf2010}; es costoso si se invocara en cada cara, de ahí las aproximaciones. En el núcleo C++ unidimensional del proyecto el exacto no participa en el avance temporal (inspección del repositorio).

Según \textcite{toro2009}, HLL usa dos ondas y promedia el estado intermedio, lo que difumina el contacto; HLLC introduce $S^{*}$ y restaura esa onda; HLLE se distingue por la estimación de velocidades de Einfeldt, no por tener tres ondas. El código C++ actual implementa HLL (estimaciones tipo Davis) y HLLC, \textbf{no} HLLE. Comparar HLL con HLLC en un mismo problema es coherente con el prototipo; incluir HLLE o WENO es una extensión sujeta a decisión del director. \textcite{liska2003} muestran que el ordenamiento de esquemas en Euler 1D/2D depende del caso: no hay un único solver universalmente superior.

\section{Software existente}

Los criterios de comparación se derivan de los objetivos del plan (intercambiar un componente, reproducir un experimento, hacer visible el método), no del número de modelos de turbulencia.

El uso de volúmenes finitos y solvers de Riemann para las ecuaciones de Euler no es meramente teórico; existen implementaciones concretas. En \textit{Clawpack: Building an Open Source Ecosystem for Solving Hyperbolic PDEs} los solvers de Riemann comparten interfaz y un repositorio común; en PyClaw la elección puede hacerse en tiempo de ejecución e incluye HLL y HLLC para Euler unidimensional \parencite{mandli2016}. Esos resultados sustentan que la intercambiabilidad de flujos numéricos es un contrato de software ya realizado en un código maduro. Por otro lado, \textit{Comparison of Several Difference Schemes on 1D and 2D Test Problems for the Euler Equations} ejecuta varios esquemas sobre una batería común de problemas de Euler y muestra que el ordenamiento depende del caso, no de un ganador universal \parencite{liska2003}. Ello justifica un banco de pruebas en el que se varíe de a un componente, con las mismas mallas y normas.

\textcite{weller1998} y \textcite{jasak2009} sitúan OpenFOAM como código abierto de investigación e industria. No se afirma que sea imposible modificar un flujo numérico; sí que su tamaño no está pensado como laboratorio 1D de pocas semanas. \textcite{economon2016} presentan SU2 como suite de simulación y diseño. \textcite{stone2020} describen Athena++ como marco AMR con MHD.

\textcite{mandli2016} documentan Clawpack: interfaz común de solvers de Riemann; en Fortran la selección es típica en compilación y en PyClaw en tiempo de ejecución. La documentación oficial de Clawpack lista solvers Python de Euler HLL y HLLC. Por tanto \textbf{no} es cierto que no exista software abierto donde se intercambie HLL/HLLC. El hueco no es esa inexistencia.

\textcite{wilson2014} recomiendan versionado y pruebas en software científico; \textcite{ramakrishnan2017}, usabilidad; \textcite{mesnard2017}, que reproducir CFD es difícil aun con código público. \textcite{oberkampf2010} distinguen verificación (¿se resuelven bien las ecuaciones discretas?) de validación (¿el modelo representa la física medida?).

\begin{table}[htbp]
\centering
\small
\begin{tabular}{p{3.2cm}cccc}
\hline
Criterio & OpenFOAM/SU2 & Clawpack & Athena++ & Prototipo 1D de este TG \\
\hline
Euler FVM & Sí & Sí & Sí & Sí (1D) \\
Cambiar Riemann & Costoso & Sí (PyClaw) & Experto & Sí (namelist) \\
Código compacto & No & Medio & No & Objetivo de diseño \\
Trazar ecuación--código en GUI & No & Docs, no GUI UIS & No & Prototipo \\
\hline
\end{tabular}
\caption{Comparación cualitativa respecto de los objetivos del trabajo de grado, no respecto de madurez industrial.}
\end{table}

\section{Síntesis y hueco}

Los métodos están maduros \parencite{toro2009,leveque2002}. El software maduro existe y Clawpack ya permite cambiar Riemann \parencite{mandli2016}. El proyecto \textbf{propone explorar} un artefacto compacto en C++, con configuración explícita, regresión y (si se conserva el laboratorio) el recorrido ecuación--algoritmo--implementación--gráfica, para formación en Ingeniería de Sistemas, sin pretender sustituir OpenFOAM, SU2 ni Clawpack. La contribución defendible es hacer observable el efecto de combinar reconstrucción y flujo numérico en problemas de Riemann estándar, no «inventar el primer código intercambiable».

WENO, HLLE, Euler 2D en el mismo binario y JSON/YAML se tratan como \textbf{propuestos} hasta que el director los confirme. El laboratorio gráfico es prototipo.

\chapter{Metodología}

\section{Tipo de investigación}

El trabajo es investigación aplicada con desarrollo de un artefacto de software científico. La unidad de análisis es el \textit{framework} (núcleo C++ y, en el prototipo, el laboratorio Python). No se declara un estudio de usuarios ni Design Science Research en sentido institucional completo, salvo que el Comité lo exija: no hay encuesta ni resultados de aprendizaje que reportar.

\section{Ciclo de desarrollo}

El desarrollo es incremental y verificado. El núcleo unidimensional ya existe; las fases de implementación 1D se leen como mantenimiento, extensión controlada y documentación, no como reescritura. El ciclo no autoriza alterar HLLC para acomodar la interfaz gráfica.

Las fases son: (0) inventario de lo implementado frente a lo propuesto; (1) literatura y requisitos; (2) arquitectura alineada con \texttt{include/euler1d}; (3) núcleo numérico (regresión en verde); (4) configuración, CLI y prototipo de GUI sin duplicar el solver; (5) verificación y regresión; (6) inspección educativa, sin inventar usuarios; (7) manuales y prueba de reproducibilidad desde el repositorio.

El contrato de configuración vigente es el namelist \texttt{\&euler} y los archivos \texttt{.xl}. JSON o YAML, si se adoptan, son extensión.

\section{Verificación y validación}

Se adopta el vocabulario de \textcite{oberkampf2010}. La comparación con el Riemann exacto unidimensional es \textbf{verificación de código y de solución}, no validación física experimental. No hay campaña de laboratorio. Los casos 2D de la literatura solo aplican si existe el solver 2D.

La suite de tests y la comparación con el Fortran de referencia del núcleo 1D no se modifican como parte de la redacción académica. Los umbrales numéricos nuevos (normas L1/L2 oficiales del TG) quedan provisionales hasta aprobación del director.

\section{Evaluación educativa}

Si se mantiene el objetivo de documentación pedagógica, la evidencia proporcionada a un trabajo de grado es una lista de comprobación (localizar la ficha de un método, ejecutar un caso documentado, ver el archivo fuente citado). Un estudio con participantes requiere diseño ético y no se finge.

\section{Riesgos}

El principal riesgo es un alcance (WENO, HLLE, 2D, JSON, estudio de aula) mayor que el calendario. Mitigación: decisión explícita del director. Riesgo científico: que la GUI implemente otro HLLC; se prohíbe. Riesgo de lenguaje: llamar validación al tubo de Sod.
